\section{Scalar matching at all crossings}
\label{n:sec-matching-growth}

The invariant graph expresses the high angular component through the low
coordinates, reducing exact continuation to a finite-dimensional system.
It has not yet produced a line that follows each
continued reference label.  We now choose all crossing controls
simultaneously so that those lines are exactly invariant. The next section
will turn their compatible traces into entire harmonic functions.

Proposition~\ref{n:prop-angular-program} supplies isolated, linearly
split doubles with finitely repeated crossing types and a positive minimum
phase separation. Recall that $N=p-B_L$ is the band size, and set
$P_{\mathrm{lab}}=NS$, the common eigenvalue period. The continued frequencies and their means are
\begin{equation}
 \Theta_i(s)=\mathfrak d(a^{-2}\lambda_i(s)),
 \label{n:eq-continued-frequency}
\end{equation}
\begin{equation}
 m_i=\frac1{P_{\mathrm{lab}}}\int_0^{P_{\mathrm{lab}}}\Theta_i(s)\,ds<k
 \qquad(1\le i\le p).
 \label{n:eq-subcritical-means}
\end{equation}
The functions $\Theta_i$ are Lipschitz; their strict mean bound was
fixed before the radial construction. Each pair either never crosses and
has a uniform gap, or its crossings repeat with bounded phase gaps.
The constant initial phase interval separates the start from the first
crossing.

Use the fixed bump choice in which each support is contained in a phase
neighborhood where the selected coefficient in
Proposition~\ref{n:prop-outward-crossing} has one nonzero sign, and each
bump equals one on a smaller fixed neighborhood of its crossing.  The
supports, plateaux, and directions are all fixed before \(J\).  For the
fixed \(k\) and \(p\), Proposition~\ref{n:prop-exact-graphs} gives
\begin{equation}
 A_j=\operatorname{diag}(\alpha_1(j),\ldots,\alpha_p(j))
          +\mathcal E_j,\qquad
 \|\mathcal E_j\|\leq C\delta_j,
 \quad
 \alpha_i(j)=e^{\mathfrak d(b(j)^{-2}\lambda_i(s(j)))}.
 \label{n:eq-matching-matrix}
\end{equation}
The graph maps and these matrices are continuous in the full product of
controls.  At a crossing \(\nu\) of \(i\) and \(h\), on its plateau,
\begin{equation}
 (\mathcal E_j)_{hi}
 =z_\nu\eta_\nu^{3/4}m_{hi}(s(j),b(j))+o(\delta_j),
 \qquad |m_{hi}(s(j),b(j))|\geq c_0>0.
 \label{n:eq-matching-offdiagonal}
\end{equation}
The sign is fixed there, and the error is uniform with every other control
free.  The analogous formula for the reversed entry is available when that
ordering requires the match.  The common radial factor \(b\) is independent
of the controls.  All constants from now on may depend on the fixed finite
program and \(p\); increasing \(J\) never changes those data.
We denote auxiliary positive lower-bound constants by \(c_*\),
\(c_*'\), and \(c_*''\), allowing their values to change between estimates.

\subsection{One scalar obstruction at each crossing}

Write \(\mathbf e_i\) for the \(i\)-th coordinate vector of \(\mathbb R^p\);
the functions \(e_i(s(j))\) are recovered through the frame map \(\iota_j\).
For each \(i\), seek a column
\[
 v_i(j)=\mathbf e_i+\sum_{h\neq i}u_{hi}(j)\mathbf e_h.
\]
For the moment all column coordinates and controls are independent
variables in the boxes
\begin{equation}
 |u_{hi}(j)|\leq\eta(j)^{1/8}\quad(h\neq i,\ j\geq J),
 \qquad |z_\nu|\leq1.
 \label{n:eq-matching-boxes}
\end{equation}
Since \(A_j=\operatorname{diag}(\alpha_i(j))+O(\delta_j)\),
the coordinate \((A_jv_i(j))_i=\alpha_i(j)+O(\delta_j)\) is
positive on all these boxes for late \(J\). We may therefore divide by it.
The line spanned by \(v_i(j)\) is carried by \(A_j\) to the line spanned by \(v_i(j+1)\)
exactly when
\begin{equation}
 u_{hi}(j+1)=r_{hi}(j)u_{hi}(j)+F_{hi}(j),\qquad
 r_{hi}(j)=\frac{\alpha_h(j)}{\alpha_i(j)},
 \label{n:eq-matching-recursion}
\end{equation}
where
\begin{equation}
 F_{hi}(j)
 =\frac{(A_jv_i(j))_h}{(A_jv_i(j))_i}
       -r_{hi}(j)u_{hi}(j).
 \label{n:eq-matching-force}
\end{equation}
Here is the force explicitly.  Suppress the index \(j\) in the next display
and write \(\mathcal E_{h\ell}=(\mathcal E_j)_{h\ell}\).  Direct expansion
of \eqref{n:eq-matching-force} gives
\begin{equation}
 F_{hi}=
 \frac{\displaystyle
   \mathcal E_{hi}+\sum_{\ell\neq i}\mathcal E_{h\ell}u_{\ell i}
   -r_{hi}u_{hi}
       \left(\mathcal E_{ii}
             +\sum_{\ell\neq i}\mathcal E_{i\ell}u_{\ell i}\right)}
 {\displaystyle
   \alpha_i+\mathcal E_{ii}
             +\sum_{\ell\neq i}\mathcal E_{i\ell}u_{\ell i}} .
 \label{n:eq-expanded-force}
\end{equation}
The \(\alpha_i\)'s and their ratios stay in a fixed compact
positive interval because \(p\) is fixed.  Thus
\begin{equation}
 |F_{hi}(j)|\leq C_1\delta_j
 \label{n:eq-force-bound}
\end{equation}
uniformly on the whole product.  This estimate uses the exact cancellation
of the diagonal reference matrix in the force.

At a crossing of this pair, take \(z_\nu=1\) or \(-1\).  On its plateau,
\eqref{n:eq-expanded-force} and \eqref{n:eq-matching-offdiagonal} yield
\begin{equation}
 F_{hi}(j)=
 z_\nu\eta_\nu^{3/4}
     \frac{m_{hi}(s(j),b(j))}{\alpha_i(j)}
       +o(\delta_j).
 \label{n:eq-force-sign}
\end{equation}
Every numerator term involving a column coordinate is
\(O(\delta_j\eta(j)^{1/8})\); replacing the denominator by \(\alpha_i(j)\)
costs \(O(\delta_j^2)\).  Both are \(o(\delta_j)\).  On a fixed plateau
\(\eta(j)/\eta_\nu\to1\) uniformly, so the leading term in
\eqref{n:eq-force-sign} has a strict constant sign and magnitude at least
\(c_1\eta_\nu^{3/4}\), after one late choice of \(J\).  This holds with
all other controls and all column variables free in their boxes.

The sign of the scalar ratio is determined solely by the reference
eigenvalues:
\begin{equation}
 \operatorname{sign}\log r_{hi}(j)
 =\operatorname{sign}\bigl(\lambda_h(s(j))-\lambda_i(s(j))\bigr).
 \label{n:eq-ratio-sign}
\end{equation}
Indeed the same positive factor \(b(j)^{-2}\) enters both frequencies and
\(\mathfrak d\) is strictly increasing.  The zeros are simple in phase.
For every crossing of this ordered pair, cut the integer mesh at
\[
 N_\nu=\lceil t_\nu\rceil.
\]
A step \(j<N_\nu\) begins before the crossing; a step \(j\geq N_\nu\)
begins at or after it.  In particular the step \(N_\nu-1\) uses the
pre-crossing ratio even if the crossing lies inside that step.  If \(t_\nu\)
is an integer, the ratio of the single step beginning there is one.  The
estimates below allow this neutral step.

When \(r_{hi}<1\), \eqref{n:eq-matching-recursion} is stable forward in
\(j\); when \(r_{hi}>1\), it is stable backward.  At a cut \(N=N_\nu\),
the two sign changes therefore have different roles:
\[
 \begin{array}{c@{\qquad}c@{\qquad}l}
  \text{sign change of }\log r_{hi}
    &\text{stable directions}&\text{condition at the cut}\\[2pt]
  +\ \longrightarrow\ -
    &\longleftarrow\ N\ \longrightarrow
    &\text{prescribe }u_{hi}(N)=0,\\[2pt]
  -\ \longrightarrow\ +
    &\longrightarrow\ N\ \longleftarrow
    &\text{match the incoming values}.
 \end{array}
\]
Reversing the ordered pair replaces \(r_{hi}\) by
\(r_{ih}=r_{hi}^{-1}\), interchanging the two rows.  Thus a double crossing
creates exactly one scalar matching condition, although both columns pass
through it.

Call the first type a \emph{zero cut}.  At a negative-to-positive cut
\(N\), two stable propagations arrive.  If
\(a<N<b\) are the adjacent zero cuts, the values
coming forward from \(a\) and backward from \(b\) are respectively
\begin{align}
 u_N^-&=\sum_{q=a}^{N-1}F_{hi}(q)
           \prod_{\ell=q+1}^{N-1}r_{hi}(\ell),
 \label{n:eq-forward-trial}\\
 u_N^+&=-\sum_{q=N}^{b-1}F_{hi}(q)
           \prod_{\ell=N}^{q}r_{hi}(\ell)^{-1}.
 \label{n:eq-backward-trial}
\end{align}
If the initial interval is forward stable, use \(a=J\) and trial value
zero there.  If it is backward stable, propagate backward from its first
zero cut.  A pair that never crosses has a uniform strict ratio gap:
use zero at \(J\) in the forward-stable case, and the convergent backward
sum from infinity in the backward-stable case.  A pair that does cross has
infinitely many alternating cuts with bounded phase gaps, so its other
propagations are finite.

The obstruction at a negative-to-positive cut is the mismatch
\begin{equation}
 \begin{split}
 D_\nu=u_N^--u_N^+
  ={}&\sum_{q=a}^{N-1}W^-_{q,N}F_{hi}(q)
       +\sum_{q=N}^{b-1}W^+_{N,q}F_{hi}(q),\\
 W^-_{q,N}={}&\prod_{\ell=q+1}^{N-1}r_{hi}(\ell)>0,\qquad
 W^+_{N,q}=\prod_{\ell=N}^{q}r_{hi}(\ell)^{-1}>0.
 \end{split}
 \label{n:eq-matching-mismatch}
\end{equation}
Both weight families are positive because the backward formula has a minus
sign.  This is the scalar matching condition assigned to crossing \(\nu\).

Near a simple crossing, summing the linearly vanishing gap gives Gaussian
decay of the stable products, with effective length
\(\eta_\nu^{-1/2}\).  Forcing of size \(\eta_\nu^{3/4}\) therefore gives
the scale \(\eta_\nu^{-1/2}\eta_\nu^{3/4}=\eta_\nu^{1/4}\), which fits
inside the coordinate box of radius \(\eta_\nu^{1/8}\).  The next lemma
proves these estimates uniformly and shows that the weighted sum retains
the sign imposed by the control.

\begin{lemma}[Stable weights and the sign of a mismatch]
\label{n:lem-matching-weights}
For all sufficiently late \(J\), every trial coordinate defined above
satisfies, uniformly on the full product of boxes,
\begin{equation}
 |u_{hi}^{\mathrm{trial}}(j)|\leq C\eta(j)^{1/4}
       <\eta(j)^{1/8}.
 \label{n:eq-trial-small}
\end{equation}
Orient the coordinate \(z_\nu\) once according to the sign of its selected
coefficient.  Each mismatch then satisfies
\begin{equation}
 D_\nu|_{z_\nu=1}\geq c_*\eta_\nu^{1/4}>0,\qquad
 D_\nu|_{z_\nu=-1}\leq-c_*\eta_\nu^{1/4}<0,
 \label{n:eq-mismatch-endpoint-signs}
\end{equation}
independently of every other control and column coordinate.
\end{lemma}

\begin{proof}
For a fixed ordered pair define the reference logarithmic ratio at phase
\(s\) and radial factor \(b\) by
\[
 \ell_{hi}(s,b)
 =\mathfrak d(b^{-2}\lambda_h(s))
       -\mathfrak d(b^{-2}\lambda_i(s)).
\]
On a sign interval \([s_a,s_b]\) bounded by successive zeros, simple
splitting and compactness of the finitely repeated types give
\begin{equation}
 |\ell_{hi}(s,b)|
 \geq c_*\min\{s-s_a,s_b-s,1\},
 \label{n:eq-ratio-lower}
\end{equation}
uniformly for \(b\) near \(a\).  Near either endpoint there is also the
upper bound \(C\) times the distance to that endpoint.  The lower bound
follows near a zero from its nonzero phase derivative and away from all
zeros from compactness.  The common radial factor does not alter the zeros,
and \(\mathfrak d'\) has positive upper and lower bounds on the finite
spectral range.  The phase lengths of these intervals have positive lower
and finite upper bounds.  Initial intervals are merely truncated subsets
of a sign interval for the periodic reference data, so the same estimates
apply to them after periodically extending the reference data for this
comparison.

Over a bounded phase interval, the elapsed time is \(O(t^{3/4})=o(t)\).
Thus the speed at every point of one such interval is comparable to a
single value \(\eta_*\), with relative error tending to zero as \(J\)
increases.  Consecutive mesh phases are separated by comparable multiples
of \(\eta_*\).  Consider a stable product over \(m\) consecutive steps
inside the interval.  Apart from a bounded number of endpoint rounding
steps, at least a fixed fraction of those phases have distance at least
\(c_*\eta_*m\) from both endpoints.  To see this, phases within distance
\(d\) of the two endpoints number at most \(C(d/\eta_*+1)\); choose
\(d=c_*'\eta_*m\) with \(c_*'\) small.  Since \(\eta_*m\) is bounded by the
fixed maximum phase length, \(\min\{d,1\}\geq c_*''\eta_*m\).
Summing \eqref{n:eq-ratio-lower} along the product therefore gives
\[
 \sum |\log r_{hi}|\geq c_*\eta_*m^2-C.
\]
The sign in a stable direction makes the product the exponential of the
negative of this sum.  The one-step ceiling convention, including a
possible neutral step, only changes the constant.  Every stable weight
therefore obeys
\begin{equation}
 W(m)\leq C e^{-c_*\eta_*m^2},\qquad
 \sum_{m\geq0}W(m)\leq C\eta_*^{-1/2}.
 \label{n:eq-gaussian-green}
\end{equation}
The second estimate follows by comparison with the Gaussian integral.

On a switching interval, \(\delta_q\) is comparable to
\(\eta_*^{3/4}\).  Combining \eqref{n:eq-force-bound} with the Green sum
in \eqref{n:eq-gaussian-green} bounds every stable trial propagation by
\(C\eta_*^{3/4}\eta_*^{-1/2}=C\eta_*^{1/4}\), which is comparable to
\(\eta(j)^{1/4}\) at the output step.
For a noncrossing pair compactness instead gives a fixed \(\rho<1\) with
either \(r_{hi}\leq\rho\) or \(r_{hi}^{-1}\leq\rho\) at every step.
The backward tail is bounded by
\(\sum_{q=j}^{\infty}C\rho^{q-j}\delta_q\leq C'\delta_j\), because
\(\delta_q\) decreases.  For the forward sum, whose initial interval can
be arbitrarily long, use
\[
 \frac{\delta_q}{\delta_j}=(j/q)^{9/16}
 \leq(1+j-q)^{9/16}\qquad(J\leq q\leq j).
\]
It follows that
\[
 \sum_{q=J}^{j-1}\rho^{j-1-q}\delta_q
 \leq\delta_j\sum_{m\geq0}\rho^m(m+2)^{9/16}
 \leq C'\delta_j.
\]
This is also bounded by \(C\eta(j)^{1/4}\).
Finally \(C\eta^{1/4}<\eta^{1/8}\) once \(J\) is late, proving
\eqref{n:eq-trial-small}.

For the sign at the crossing \(\nu\), take a fixed small \(\theta>0\) and
the mesh window
\[
 |j-t_\nu|\leq\theta\eta_\nu^{-1/2}.
\]
Its phase width is \(O(\theta\eta_\nu^{1/2})\), so the window lies inside
the fixed plateau for late \(J\).  The upper linear bound on
\(|\ell_{hi}|\) near the zero shows that an incoming weight from any
index in this window is at least \(\exp(-C\theta^2)\): the sum of its
logarithmic losses is at most \(C\eta_\nu\sum_{m\leq
\theta\eta_\nu^{-1/2}}(m+1)\).  The two incoming sides together contain
at least \(c_*\eta_\nu^{-1/2}\) such indices.  At \(z_\nu=\pm1\), all their
forces have the sign in \eqref{n:eq-force-sign} and magnitude at least
\(c_1\eta_\nu^{3/4}\).  Their contribution to
\eqref{n:eq-matching-mismatch} has that sign and magnitude at least
\(c_2\eta_\nu^{1/4}\).

Every other term on the plateau has the same sign.  Outside the fixed
plateau the phase distance to the crossing is bounded below by a positive
constant.  Its incoming weight contains a stable segment of length
\(c_*\eta_\nu^{-1}\); \eqref{n:eq-gaussian-green} bounds the weight by
\(C e^{-c_*'/\eta_\nu}\).  The two adjacent switching intervals contain at
most \(C\eta_\nu^{-1}\) indices, and their forces are
\(O(\eta_\nu^{3/4})\).  Their total contribution outside the plateau is
therefore at most
\[
 C\eta_\nu^{-1/4}e^{-c_*'/\eta_\nu}
   =o(\eta_\nu^{1/4}).
\]
For a truncated initial interval, separation of the start from the first
crossing places the central window after \(J\) for late starts.  Deleting
earlier terms does not change its signed lower bound, and the retained
outside-plateau terms satisfy the same tail estimate.  The sign from the
central window dominates.  Reverse the orientation of \(z_\nu\) when
needed so that its positive endpoint gives positive mismatch.  Uniformity
of \eqref{n:eq-force-sign} over every other variable yields
\eqref{n:eq-mismatch-endpoint-signs}.
\end{proof}

\subsection{Simultaneous matching}

Consider the countable product \(\mathcal X\) of the intervals in
\eqref{n:eq-matching-boxes}, with its product topology.  Given a point of
\(\mathcal X\), compute the forces from \eqref{n:eq-matching-force}, and
then compute every stable trial propagation.  At a cut with two incoming
values, use the forward value as the output \(u\)-coordinate.  For each
control use the output
\begin{equation}
 z_\nu\longmapsto
 \Pi_{[-1,1]}(z_\nu-D_\nu),
 \label{n:eq-control-clamp}
\end{equation}
where \(\Pi_{[-1,1]}\) is interval projection.  Lemma
\ref{n:lem-matching-weights} puts each trial coordinate strictly inside
its box, and the projection puts each control inside its interval.  These
rules define a self-map \(\Phi:\mathcal X\to\mathcal X\).

This map is continuous in product topology.  For a fixed \(j\), \(A_j\)
is continuous in the control product by Proposition~\ref{n:prop-exact-graphs},
and the force at that step depends on only finitely many column variables
through a denominator bounded away from zero.  Hence each force is
continuous.  A finite trial propagation or mismatch is a finite expression
in such forces.  In the only infinite case, a backward propagation for a
noncrossing pair, the geometric tail is uniformly summable on
\(\mathcal X\).  Its output is therefore a uniform limit of continuous
finite sums.  Each output coordinate of \(\Phi\) is continuous, as
required for the product topology.

We can obtain a fixed point using only finite-dimensional Brouwer and a
diagonal limit.  Enumerate the coordinates of \(\mathcal X\).  Freeze those
after the first \(N\) at their box centers and apply Brouwer to the first
\(N\) output coordinates; this gives a point \(x^{[N]}\in\mathcal X\)
whose first \(N\) coordinates satisfy their fixed-point equations.
By repeated subsequence extraction in the compact coordinate intervals,
choose a subsequence converging in every coordinate to \(x\in\mathcal X\).
For any fixed coordinate, its fixed-point equation holds for all
sufficiently large members of the subsequence.  Continuity of that output
coordinate passes the equation to the limit.  Thus \(\Phi(x)=x\).

At a fixed point, no control can equal \(1\): by
\eqref{n:eq-mismatch-endpoint-signs}, the argument \(1-D_\nu\) of
\eqref{n:eq-control-clamp} is strictly less than \(1\), so its projection
is less than \(1\).  Likewise a control equal to \(-1\) would have projected
image greater than \(-1\).  Every fixed control is interior, where
\(z_\nu=\Pi_{[-1,1]}(z_\nu-D_\nu)\) forces \(D_\nu=0\).
At each two-valued cut the forward trial value consequently equals the
backward one.  Away from cuts the trial propagations already obey
\eqref{n:eq-matching-recursion}; at zero cuts they share the prescribed
zero.  The fixed point therefore satisfies every scalar recurrence
exactly at every step.

Fix this control sequence and write \(g=g_z\).  Define the positive line
multiplier
\[
 \sigma_i(j)=(A_jv_i(j))_i
             =\alpha_i(j)+O(\delta_j)>0.
\]
The recurrences and the normalization of the \(i\)-th coordinate give
\begin{equation}
 A_jv_i(j)=\sigma_i(j)v_i(j+1).
 \label{n:eq-exact-lines}
\end{equation}
The \(p\) columns form a basis for late \(J\): their matrix is the identity
plus a matrix of norm \(O(p\eta(j)^{1/8})<1\).  This is one basis for each
step of one metric; the control sequence is common to all labels.
